% Benki ya mazoezi — sura ya 4
% Sura hubainishwa na jina la faili hili.

\begin{exo}[Kukokotoa integrali ya umbo tofauti kwenye njia]{integrale-forme-differentielle}
\keywords{umbo tofauti, diferenshali kamili, kigezo cha Schwarz, integrali ya mstari}

Fikiria maumbo mawili tofauti
\[
\omega_1 = y\,\mathrm{d}x + x\,\mathrm{d}y,
\qquad
\omega_2 = y\,\mathrm{d}x,
\]
na njia tatu zinazounganisha asili $O(0,0)$ na nukta $M(1,1)$: njia
$\gamma_1$ ni kipande cha mlalo $O\to(1,0)$ kikifuatwa na kipande cha wima
$(1,0)\to M$; njia $\gamma_2$ ni kipande cha wima $O\to(0,1)$ kikifuatwa
na kipande cha mlalo $(0,1)\to M$; na njia $\gamma_3$ ni kipande cha mstari
kinachounganisha $O$ moja kwa moja na $M$.

\begin{enumerate}
\item Onyesha kuwa $\omega_1$ ni kamili kwa kupata funksioni $f$ ambayo $\omega_1=\mathrm{d}f$.
\item Kokotoa $\int_{\gamma_1}\omega_1$ na $\int_{\gamma_2}\omega_1$ kwa kuweka parameta kwa kila kipande. Pata tena jibu bila kukokotoa.
\item Tumia kigezo cha Schwarz kwa $\omega_2$. Tunaweza kuhitimisha nini?
\item Kokotoa $\int_{\gamma_1}\omega_2$, $\int_{\gamma_2}\omega_2$ na $\int_{\gamma_3}\omega_2$. Toa maoni.
\end{enumerate}

\begin{indication}
Kwenye kipande cha mlalo $\mathrm{d}y=0$, na kwenye kipande cha wima
$\mathrm{d}x=0$: kila integrali hupunguzwa kuwa integrali rahisi.
\end{indication}

\begin{solution}
\textbf{Swali la 1.} $f(x,y)=xy$ inafaa, kwa sababu $\partial f/\partial x=y$ na $\partial f/\partial y=x$.

\textbf{Swali la 2.} Kwa mkunjo uliowekewa parameta kwa
$t\mapsto\bigl(x(t),y(t)\bigr)$, tuna
\[
\int_\gamma\omega_1
= \int \left[y(t)x'(t)+x(t)y'(t)\right]\,\mathrm{d}t.
\]

Njia $\gamma_1$ ni muungano wa vipande viwili. Kwenye cha kwanza,
$x(t)=t$ na $y(t)=0$, huku $t\in[0,1]$; hivyo $\mathrm{d}x=\mathrm{d}t$
na $\mathrm{d}y=0$. Kwenye cha pili, $x(t)=1$ na $y(t)=t$, bado
$t\in[0,1]$; hivyo $\mathrm{d}x=0$ na $\mathrm{d}y=\mathrm{d}t$. Kwa hiyo,
\[
\begin{aligned}
I_1
= \int_{\gamma_1}\omega_1
&= \int_0^1\left(0\times 1+t\times 0\right)\,\mathrm{d}t
 + \int_0^1\left(t\times 0+1\times 1\right)\,\mathrm{d}t \\
&= 0+\int_0^1\mathrm{d}t = 1.
\end{aligned}
\]

Kwa $\gamma_2$, kipande cha kwanza kina $x(t)=0$, $y(t)=t$, kisha cha pili
kina $x(t)=t$, $y(t)=1$, huku $t\in[0,1]$ katika hali zote mbili. Hivyo
\[
\begin{aligned}
I_2
= \int_{\gamma_2}\omega_1
&= \int_0^1\left(t\times 0+0\times 1\right)\,\mathrm{d}t
 + \int_0^1\left(1\times 1+t\times 0\right)\,\mathrm{d}t \\
&= 0+\int_0^1\mathrm{d}t = 1.
\end{aligned}
\]
Thamani hizi mbili ni sawa. Tungeweza kutabiri bila hesabu: kwa kuwa
$\omega_1=\mathrm{d}f$ na $f(x,y)=xy$, integrali yake hutegemea ncha pekee,
\[
\int_\gamma\omega_1=f(M)-f(O)=f(1,1)-f(0,0)=1.
\]

\textbf{Swali la 3.} Tuandike $\omega_2=A(x,y)\,\mathrm{d}x+B(x,y)\,\mathrm{d}y$.
Hapa $A(x,y)=y$ na $B(x,y)=0$. Kama umbo lingekuwa kamili, kigezo cha Schwarz kingetoa
\[
\frac{\partial A}{\partial y}=\frac{\partial B}{\partial x}.
\]
Lakini $\partial A/\partial y=1$ ilhali $\partial B/\partial x=0$.
Derivativi mseto ni tofauti: kwa hiyo $\omega_2$ si kamili. Integrali yake
inaweza kutegemea njia kati ya $O$ na $M$, kama hesabu inayofuata inavyoonyesha.

\textbf{Swali la 4.} Kwa kuwa $\omega_2=y\,\mathrm{d}x$, ni badiliko la $x$
linalofanyika wakati $y$ si sifuri pekee linaloweza kuchangia integrali.
Kwa kutumia tena uparametishaji wa awali, kwenye $\gamma_1$ tunapata
\[
\begin{aligned}
J_1
=\int_{\gamma_1}\omega_2
&=\int_0^1 0\times 1\,\mathrm{d}t
  +\int_0^1 t\times 0\,\mathrm{d}t \\
&=0.
\end{aligned}
\]
Kwenye $\gamma_2$, kipande cha kwanza kina $x(t)=0$, hivyo $\mathrm{d}x=0$;
kwenye cha pili, $x(t)=t$, $y(t)=1$ na $\mathrm{d}x=\mathrm{d}t$. Basi
\[
\begin{aligned}
J_2
=\int_{\gamma_2}\omega_2
&=\int_0^1 t\times 0\,\mathrm{d}t
  +\int_0^1 1\times 1\,\mathrm{d}t \\
&=1.
\end{aligned}
\]
Hatimaye, kipande cha ulalo $\gamma_3$ kinawekwa parameta waziwazi kwa
\[
x(t)=t,\qquad y(t)=t,\qquad t\in[0,1],
\]
hivyo $\mathrm{d}x=\mathrm{d}t$ na $\mathrm{d}y=\mathrm{d}t$. Tunapata
\[
J_3
=\int_{\gamma_3}\omega_2
=\int_0^1 y(t)x'(t)\,\mathrm{d}t
=\int_0^1 t\,\mathrm{d}t
=\frac12.
\]
Njia zote tatu zina ncha zilezile lakini zinatoa thamani tatu tofauti:
$J_1=0$, $J_2=1$ na $J_3=1/2$. Huu ndio utegemezi wa njia unaotambulisha
umbo tofauti lisilo kamili.
\end{solution}
\end{exo}

\begin{exo}[Badiliko lilelile la nishati ya ndani, njia tatu za uhamisho]{meme-delta-u-trois-transferts}
\keywords{kanuni ya kwanza, nishati ya ndani, joto, kazi ya kimakanika, kazi ya umeme, kalorimetria, funksioni ya hali}

Kioevu, kalorimita yake na kinzani kidogo kilichozamishwa vinaunda mfumo
uliofungwa na ulio katika mapumziko kwa kiwango cha makroskopiki. Chombo ni
kigumu na uwezo wa jumla wa joto wa mfumo unachukuliwa kuwa thabiti:
\[
C=2{,}00\ \mathrm{kJ\,K^{-1}}.
\]
Tunataka kuhamisha mfumo kutoka hali ya msawazo $A$, yenye
$T_A=20{,}0\,^\circ\mathrm{C}$, hadi hali ya msawazo $B$, yenye
$T_B=25{,}0\,^\circ\mathrm{C}$. Tunakubali kwamba
\[
U(B)-U(A)=C(T_B-T_A).
\]
Mabadiliko ya nishati ya mwendo na nishati ya uwezo ya makroskopiki ni sifuri.
Upotevu wote kwenda nje unapuuzwa.

Itifaki tatu zifuatazo zinatekelezwa kila moja peke yake kutoka hali ileile ya mwanzo $A$:
\begin{itemize}
\item[Itifaki ya 1] Mfumo unawekwa katika mgusano wa joto na chanzo cha moto bila kupewa kazi yoyote.
\item[Itifaki ya 2] Kuta zimewekewa kihami joto. Pala inakoroga kioevu kwa kuanguka kwa tungamo $m$ kwa urefu $h=5{,}00$ m. Pala na kioevu huishia katika mapumziko, na nishati yote ya kimakanika iliyopotezwa na tungamo huhamishwa kwa mfumo. Tumia $g=9{,}81\ \mathrm{m\,s^{-2}}$.
\item[Itifaki ya 3] Mfumo unapokea joto $Q_3=4{,}00$ kJ na nishati iliyobaki inapelekwa kwa kinzani kwa njia ya umeme. Nguvu ya umeme inayopokelewa ni thabiti, $\mathcal P=300$ W.
\end{itemize}

\begin{enumerate}
\item Kokotoa badiliko la nishati ya ndani $\Delta U=U(B)-U(A)$ linalofanana kwa itifaki zote tatu.
\item Kwa kila itifaki mbili za kwanza, bainisha joto $Q$ na kazi $W$ zinazopokelewa na mfumo. Katika itifaki ya pili, kokotoa tungamo $m$ kinachohitajika.
\item Kwa itifaki ya tatu, kokotoa kazi ya umeme inayopokelewa na muda wa kuipa kinzani umeme.
\item Mabadiliko yote matatu yana hali zilezile za mwanzo na mwisho. Eleza kwa nini thamani zao za $Q$ na $W$ zinaweza hata hivyo kutofautiana. Je, mfumo una “joto” au “kazi” katika hali ya mwisho $B$?
\end{enumerate}

\begin{indication}
Tumia $\Delta U=Q+W$ kwa kanuni ya ishara ya benki. Kwa pala, kazi inayopokelewa
ni sawa na upungufu $mgh$ wa nishati ya uwezo ya tungamo. Nishati ya umeme
inayovuka mpaka kupitia nyaya huhesabiwa kama kazi ya umeme.
\end{indication}

\begin{solution}
\textbf{Swali la 1.} Tofauti ya halijoto ni $T_B-T_A=5{,}0$ K. Kwa hiyo,
\[
\Delta U=C(T_B-T_A)
=2{,}00\times5{,}0
=10{,}0\ \mathrm{kJ}.
\]
Thamani hii hutegemea hali mbili za msawazo $A$ na $B$ pekee.

\textbf{Swali la 2.} Katika itifaki ya kwanza hakuna kazi inayopokelewa:
$W_1=0$. Kanuni ya kwanza basi inatoa
\[
Q_1=\Delta U=10{,}0\ \mathrm{kJ}.
\]
Joto ni chanya kwa sababu nishati inaingia kwenye mfumo.

Katika itifaki ya pili, kuta zimewekewa kihami joto, kwa hiyo $Q_2=0$.
Badiliko lote la nishati ya ndani linatokana na kazi ya kimakanika ya pala:
\[
W_2=\Delta U=10{,}0\ \mathrm{kJ}.
\]
Kwa kuwa $W_2=mgh$, tunapata
\[
m=\frac{W_2}{gh}
=\frac{10{,}0\times10^3}{9{,}81\times5{,}00}
\simeq 204\ \mathrm{kg}.
\]
Ukubwa huu mkubwa unatukumbusha kuwa hata ongezeko dogo la halijoto tayari
linalingana na nishati kubwa ya kimakanika.

\textbf{Swali la 3.} Kanuni ya kwanza inalazimisha
\[
W_3=\Delta U-Q_3=10{,}0-4{,}00=6{,}00\ \mathrm{kJ}.
\]
Kazi hii ni chanya: ugavi wa umeme unatoa nishati kwa mfumo. Kwa
$W_3=\mathcal P\,\Delta t$,
\[
\Delta t=\frac{W_3}{\mathcal P}
=\frac{6{,}00\times10^3}{300}
=20{,}0\ \mathrm{s}.
\]

\textbf{Swali la 4.} Njia tatu zinatoa mtawalia
\[
(Q_1,W_1)=(10{,}0,0)\ \mathrm{kJ},
\qquad
(Q_2,W_2)=(0,10{,}0)\ \mathrm{kJ},
\]
na
\[
(Q_3,W_3)=(4{,}00,6{,}00)\ \mathrm{kJ}.
\]
Katika kila hali, jumla $Q+W$ ni $10{,}0$ kJ na hutoa badiliko lilelile
$\Delta U$. Nishati ya ndani ni funksioni ya hali, ilhali joto na kazi
hueleza uhamisho wakati wa mabadiliko. Katika hali ya mwisho $B$, mfumo una
nishati ya ndani, lakini hauna “joto” wala “kazi” ndani yake.
\end{solution}
\end{exo}

\begin{exo}[Mbanano chini ya shinikizo la nje lisilobadilika]{compression-pression-exterieure-constante}
\seoready{true}
\keywords{kazi ya nguvu za shinikizo, shinikizo la nje lisilobadilika, mbanano, upanuzi, upanuzi kwenye ombwe}

Gesi inabanwa kutoka ujazo $V_A$ hadi ujazo $V_B<V_A$ chini ya shinikizo la nje lisilobadilika $P_0$.

\begin{enumerate}
\item Kokotoa kazi inayopokelewa na gesi.
\item Gesi inarudi kutoka $V_B$ hadi $V_A$ dhidi ya shinikizo lilelile la nje $P_0$. Kokotoa kazi inayopokelewa na uilinganishe na kazi ya mbanano.
\item Rudia upanuzi kutoka $V_B$ hadi $V_A$ wakati gesi inapanuka kwenye ombwe. Fasiri ishara zote tatu zilizopatikana.
\end{enumerate}

\begin{indication}
Tumia $W=-\int P_{\mathrm{ext}}\,\mathrm{d}V$. Lazima tukokotoe integrali ya
shinikizo la \emph{nje}, hata wakati mabadiliko si ya nusu-tuli.
\end{indication}

\begin{solution}
\textbf{Swali la 1.} Shinikizo la nje ni thabiti, kwa hiyo
\[
W_{A\to B}=-P_0(V_B-V_A)=P_0(V_A-V_B)>0.
\]
Gesi inapokea kazi wakati wa mbanano.

\textbf{Swali la 2.} Kwa njia ya kurudi,
\[
W_{B\to A}=-P_0(V_A-V_B)<0.
\]
Kazi hizi mbili zina ishara kinyume kwa sababu mabadiliko yanafanywa dhidi
ya shinikizo lilelile la nje lisilobadilika.

\textbf{Swali la 3.} Kwenye ombwe, $P_{\mathrm{ext}}=0$, kwa hiyo
\[
W_{B\to A}^{\mathrm{vide}}=0.
\]
Kwa hiyo upanuzi haumaanishi kazi hasi moja kwa moja: lazima nguvu ya nje
isukumwe kweli.
\end{solution}
\end{exo}
